\documentclass[12pt,a4paper]{article}
\usepackage{fontspec}
\setmainfont{DejaVu Serif}
\setsansfont{DejaVu Sans}
\setmonofont{DejaVu Sans Mono}[Scale=0.85]
\usepackage[margin=2.2cm]{geometry}
\usepackage{amsmath}
\usepackage{amssymb}
\usepackage{xcolor}
\usepackage{enumitem}
\usepackage{booktabs}
\usepackage{array}
\usepackage{listings}
\usepackage{hyperref}
\usepackage{titlesec}

\titleformat{\section}{\large\bfseries}{}{0pt}{}
\setlength{\parskip}{0.5em}
\setlength{\parindent}{0pt}
\sloppy
\setlength{\emergencystretch}{3em}
\hypersetup{hidelinks}

\lstset{
  basicstyle=\ttfamily\small,
  keywordstyle=\color{blue!70!black}\bfseries,
  commentstyle=\color{green!45!black},
  frame=single, framerule=0.3pt, rulecolor=\color{gray!60},
  backgroundcolor=\color{gray!6},
  columns=fullflexible, keepspaces=true,
  language=C, showstringspaces=false,
  xleftmargin=1em, xrightmargin=1em
}

\newcommand{\question}[2]{\section{Дасгал #1}\textbf{Асуулт.} #2\par\medskip\textbf{Хариулт.}\par}

\begin{document}

\begin{center}
  {\Large\bfseries Лабораторийн ажил №9}\\[0.4em]
  {\large Виртуал санах ой (Virtual Memory) — Бүлэг 10, Practice Exercises}\\[0.8em]
  F.CSM302 Үйлдлийн систем\\
  Гүйцэтгэсэн: Б.Билгүүн (B190910014)\\[0.3em]
  Сурах бичиг: A. Silberschatz, P. B. Galvin, G. Gagne, \textit{Operating System Concepts}, 10th ed., Wiley, 2018
\end{center}
\hrule
\bigskip

Виртуал санах ой нь процессын логик хаягийн орон зайг физик санах ойгоос салгаж, программыг бүхэлд нь санах ойд ачаалалгүйгээр ажиллуулах боломж олгодог. \textbf{Шаардлагаар хуудаслах} (demand paging) үед хуудсыг зөвхөн хандах үед нь санах ойд авчирна; санах ойд байхгүй хуудсанд хандвал \textbf{page fault} үүсч, чөлөөт фрэйм байхгүй бол \textbf{хуудас солих алгоритм} (FIFO, LRU, OPT, second-chance) нэг хуудсыг гаргана. Хүснэгтийн тоонуудыг \texttt{lab9/vmcalc.py} ба \texttt{lab9/page\_replace.c} хоёр тусдаа хэрэгжүүлэлтээр шалгасан.

%---------------------------------------------------------------
\question{10.1}{Page fault ямар нөхцөлд үүсдэг вэ? Page fault үүсэхэд үйлдлийн систем ямар үйлдлүүд хийдэг вэ?}

Процесс санах ойд ачаалагдаагүй хуудсанд хандахад — хуудасны хүснэгтийн бичлэгийн valid–invalid бит \emph{invalid} байх үед — хаяг хөрвүүлэх техник хангамж (MMU) үйлдлийн систем рүү trap үүсгэнэ. Үйлдлийн системийн алхмууд (§10.2.1):
\begin{enumerate}[leftmargin=1.5em]
  \item Процессын дотоод хүснэгт (ихэвчлэн PCB-д)-ээс хандалт хүчинтэй эсэхийг шалгана. Хүчингүй (хаягийн орон зайнаас гадуур) бол процессыг зогсооно; хүчинтэй ч хуудас нь санах ойд байхгүй бол үргэлжлүүлнэ.
  \item Чөлөөт фрэйм олно (чөлөөт фрэймийн жагсаалтаас; байхгүй бол хуудас солих алгоритмаар нэг хуудсыг гаргаж, өөрчлөгдсөн бол (dirty) дискэнд бичнэ).
  \item Хэрэгтэй хуудсыг хоёрдогч санах ой (swap)-аас тэр фрэйм рүү унших дискний үйлдлийг товлоно. Энэ хооронд процесс хүлээлтийн төлөвт орж, CPU өөр процесст шилжиж болно.
  \item Уншилт дуусахад процессын дотоод хүснэгт ба хуудасны хүснэгтийг шинэчилж, бичлэгийг \emph{valid} болгоно.
  \item Trap-аар тасалдсан зааврыг \textbf{дахин эхлүүлнэ} — процесс хуудас үргэлж санах ойд байсан мэт үргэлжилнэ.
\end{enumerate}

%---------------------------------------------------------------
\question{10.2}{$m$ фрэймтэй (эхэндээ бүгд хоосон) процессын хуудасны лавлагааны мөр $p$ урттай, түүнд $n$ ялгаатай хуудасны дугаар орно. Ямар ч хуудас солих алгоритмын хувьд: (a) page fault-ийн доод хязгаар хэд вэ? (b) дээд хязгаар хэд вэ?}

\begin{enumerate}[label=(\alph*), leftmargin=1.5em]
  \item \textbf{Доод хязгаар $n$.} Фрэймүүд эхэндээ хоосон тул ялгаатай хуудас бүр анх хандахдаа заавал нэг fault үүсгэнэ — ямар ч алгоритм үүнээс бага байж чадахгүй ($m \ge n$ үед яг $n$).
  \item \textbf{Дээд хязгаар $p$.} Хандалт бүр хамгийн ихдээ нэг fault үүсгэнэ — бүх хандалт fault болох нь хамгийн муу тохиолдол (жнь $m < n$ үед FIFO/LRU циклийн хандалт дээр).
\end{enumerate}

%---------------------------------------------------------------
\question{10.3}{LRU, FIFO, Optimal, Second-chance хуудас солих алгоритмуудыг page fault-ийн түвшнээр нь “муу”-гаас “төгс” хүртэл таван онооны шатлалаар эрэмбэл. Belady-ийн аномалид өртдөг ба өртдөггүйг ялга.}

\begin{center}
\begin{tabular}{cl>{\raggedright}p{5.2cm}l}
\toprule
Оноо & Алгоритм & Үнэлгээ & Belady\\
\midrule
5 & Optimal (OPT) & төгс — онолын хамгийн бага fault & өртөхгүй\\
4 & LRU & сайн — OPT-д ойр & өртөхгүй\\
3 & Second-chance & дунд — LRU-ийн ойролцоо & өртдөг\\
1–2 & FIFO & муу & өртдөг\\
\bottomrule
\end{tabular}
\end{center}
OPT ба LRU бол \textbf{стек алгоритм}: $k$ фрэймтэй үед санах ойд байх хуудсын олонлог үргэлж $k+1$ фрэймтэй үеийнхийн дэд олонлог тул фрэйм нэмэхэд fault хэзээ ч өсөхгүй. FIFO ба (бүх reference бит 1 үед FIFO болж хувирдаг) second-chance нь энэ шинжгүй тул фрэйм нэмэхэд fault өсөх тохиолдол гарна.

%---------------------------------------------------------------
\question{10.4}{CPU-ийн циклийн хугацаа 1 мкс; одоогийнхоос өөр хуудсанд хандах нэмэлт 1 мкс. Хуудас 1000 үгтэй; хуудасны төхөөрөмж нь минутад 3000 эргэдэг, секундэд 1 сая үг дамжуулдаг дугуй (drum). Нийт зааврын 1\% нь өөр хуудсанд хандсан; тэдгээрийн 80\% нь санах ойд байгаа хуудсанд; шинэ хуудас хэрэгтэй үед солигдох хуудас 50\%-д нь өөрчлөгдсөн байсан. Нэг процесс ажиллаж, дамжуулалтын үед процессор сул зогсдог гэж үзээд зааврын үр ашигтай хугацааг ол.}

\textbf{Хуудас дамжуулах хугацаа:} 3000 эрг/мин $\Rightarrow$ нэг эргэлт 20 мс, дундаж эргэлтийн хүлээлт 10 мс; дамжуулах $1000 / 10^6$ с $= 1$ мс. Нэг хуудас унших/бичих $\approx 11$ мс $= 11\,000$ мкс.

\begin{center}
\begin{tabular}{lrr}
\toprule
Тохиолдол & Магадлал & Хугацаа (мкс)\\
\midrule
Одоогийн хуудсанд & 0.99 & 1\\
Өөр хуудас, санах ойд & $0.01 \times 0.8 = 0.008$ & 2\\
Page fault (унших) & $0.01 \times 0.2 = 0.002$ & $2 + 11\,000$\\
\quad нэмэлт: өөрчлөгдсөн хуудсыг бичих & $0.002 \times 0.5 = 0.001$ & $11\,000$\\
\bottomrule
\end{tabular}
\end{center}
\[
T = 0.99 \cdot 1 + 0.008 \cdot 2 + 0.002 \cdot 11\,002 + 0.001 \cdot 11\,000 = 0.99 + 0.016 + 22.004 + 11 \approx \mathbf{34\ мкс}.
\]
Өөрөөр хэлбэл 1000 зааврын 2-хон page fault зааврын дундаж хугацааг 34 дахин удаашруулна — page fault-ийн түвшинг маш бага байлгах шаардлагатайг харуулж байна.

%---------------------------------------------------------------
\question{10.5}{12 бит виртуал ба физик хаягтай, 256 байт хуудастай системийн хуудасны хүснэгт өгөгдсөн (хуудас 0–9 → –, 2, C, A, –, 4, 3, –, B, 0; “–” нь санах ойд байхгүй). Чөлөөт фрэймийн жагсаалт D, E, F (D толгойд). 9EF, 111, 700, 0FF (16-тын) виртуал хаягийг физик хаяг руу хөрвүүл.}

256 $= 2^8$ байт хуудас → доод 8 бит (сүүлийн 2 hex орон) оффсет, дээд 4 бит (эхний hex орон) хуудасны дугаар. Физик хаяг $=$ фрэйм$\,\|\,$оффсет.
\begin{center}
\begin{tabular}{cccccl}
\toprule
Виртуал & Хуудас & Оффсет & Фрэйм & Физик & Тайлбар\\
\midrule
9EF & 9 & EF & 0 & \textbf{0EF} & санах ойд байна\\
111 & 1 & 11 & 2 & \textbf{211} & санах ойд байна\\
700 & 7 & 00 & D & \textbf{D00} & page fault → чөлөөт D\\
0FF & 0 & FF & E & \textbf{EFF} & page fault → чөлөөт E\\
\bottomrule
\end{tabular}
\end{center}

%---------------------------------------------------------------
\question{10.6}{Шаардлагаар хуудаслах (demand paging)-ыг дэмжихэд шаардлагатай техник хангамжийн функцуудыг хэлэлц.}

\begin{itemize}[leftmargin=1.5em]
  \item \textbf{Valid–invalid (эсвэл protection) бит бүхий хуудасны хүснэгт} — хуудас санах ойд байгаа эсэхийг тэмдэглэж, invalid бит дээр хандахад page fault trap үүсгэх MMU.
  \item \textbf{Хоёрдогч санах ой (swap space)} — санах ойд байхгүй хуудсуудыг хадгалах өндөр хурдны төхөөрөмж.
  \item \textbf{Зааврыг дахин эхлүүлэх чадвар} — page fault-аар тасалдсан зааврыг (олон операндтай, санах ойн блок зөөдөг зааврыг ч) яг тэр төлөвөөс нь, хагас хийгдсэн өөрчлөлтийг буцааж эсвэл урьдчилан шалгаж дахин ажиллуулах.
  \item Хуудас солих алгоритмд хэрэгтэй \textbf{reference бит} ба \textbf{dirty (modify) бит} — сүүлд хэрэглэгдсэн эсэх, бичигдсэн эсэхийг техник хангамж автоматаар тэмдэглэнэ (dirty биш хуудсыг дискэнд бичихгүйгээр гаргана).
  \item Хаяг хөрвүүлэлтийг хурдасгах \textbf{TLB}.
\end{itemize}

%---------------------------------------------------------------
\question{10.7}{\texttt{int A[][] = new int[100][100];} хоёр хэмжээст массивын A[0][0] нь 200 хэмжээтэй хуудас бүхий санах ойн 200-р байрлалд байна. Массивыг боловсруулах жижиг процесс 0-р хуудсанд (0–199) байрлана. Гурван фрэймтэй (1-р фрэйм процесст, бусад хоёр хоосон), LRU солилттой үед дараах массив эхлүүлэх давталт хэдэн page fault үүсгэх вэ?\\ (a) \texttt{for j \{ for i \{ A[i][j] = 0; \}\}} \quad (b) \texttt{for i \{ for j \{ A[i][j] = 0; \}\}}}

Массив мөрөөр (row-major) хадгалагдана: нэг мөр 100 үг, нэг хуудас 200 үг тул \textbf{нэг хуудсанд 2 мөр}, массив бүхэлдээ 50 хуудас. Өгөгдөлд ердөө 2 фрэйм байна.
\begin{enumerate}[label=(\alph*), leftmargin=1.5em]
  \item Гадна давталт багана $j$, дотор давталт мөр $i$: нэг баганыг бөглөхөд 100 мөр буюу бүх 50 хуудсаар дамжина (мөр бүр хоёр дахь дараалсан хандалт нь ижил хуудас). 2 фрэйм 50 хуудсыг барьж чадахгүй тул багана бүрт 50 fault: $100 \times 50 = \mathbf{5000}$ page fault.
  \item Гадна давталт мөр $i$: санах ойд дараалан хандана, хуудас бүр нэг удаа л ачаалагдана: $\mathbf{50}$ page fault.
\end{enumerate}
Дүгнэлт: массивыг санах ойд хадгалагдсан дарааллаар нь (locality) явуулах нь page fault-ийг 100 дахин багасгана.

%---------------------------------------------------------------
\question{10.8}{Хуудасны лавлагааны мөр: 1, 2, 3, 4, 2, 1, 5, 6, 2, 1, 2, 3, 7, 6, 3, 2, 1, 2, 3, 6. Нэгээс долоон фрэймтэй үед LRU, FIFO, Optimal солилтод хэдэн page fault үүсэх вэ? (Фрэймүүд эхэндээ хоосон.)}

\begin{center}
\begin{tabular}{lrrrrrrr}
\toprule
Фрэймийн тоо & 1 & 2 & 3 & 4 & 5 & 6 & 7\\
\midrule
LRU     & 20 & 18 & 15 & 10 & 8 & 7 & 7\\
FIFO    & 20 & 18 & 16 & 14 & 10 & 10 & 7\\
Optimal & 20 & 15 & 11 & 8 & 7 & 7 & 7\\
\bottomrule
\end{tabular}
\end{center}
Мөрөнд 7 ялгаатай хуудас (1–7) байгаа тул 7 фрэймтэй үед бүх алгоритм зөвхөн эхний ачаалалтын 7 fault-тай; 1 фрэймтэй үед дараалсан ижил хуудас байхгүй тул хандалт бүр fault (20). Бүх баганад OPT $\le$ LRU $\le$ FIFO. Жишээ болгон LRU, 3 фрэйм: \texttt{1 2 3 4 2 1 5 6 2 1 2 3 7 6 3 2 1 2 3 6} → дахин хандалтаар fault болохгүй байрлалууд нь 5-р (2), 11-р (2), 15-р (3), 18-р (2), 19-р (3) — $20 - 5 = 15$.

%---------------------------------------------------------------
\question{10.9}{Хуудасны лавлагааны мөр: 7, 2, 3, 1, 2, 5, 3, 4, 6, 7, 7, 1, 0, 5, 4, 6, 2, 3, 0, 1. Гурван фрэймтэй шаардлагаар хуудаслах үед LRU, FIFO, Optimal солилтод хэдэн page fault үүсэх вэ?}

\begin{center}
\begin{tabular}{lc}
\toprule
Алгоритм & Page fault (3 фрэйм)\\
\midrule
LRU & \textbf{18}\\
FIFO & \textbf{17}\\
Optimal & \textbf{13}\\
\bottomrule
\end{tabular}
\end{center}
Энэ мөрөнд LRU нь FIFO-оос нэгээр \emph{муу}: хандалтууд бараг давтагдахгүй, locality муу тул “сүүлд хэрэглэгдсэн” мэдээлэл ирээдүйг сайн таамаглахгүй. LRU нь OPT-ийн ойролцоолол боловч бүх мөрөнд FIFO-оос сайн байх баталгаагүй.

%---------------------------------------------------------------
\question{10.10}{Reference бит шаарддаг хуудаслах алгоритм (second-chance, working-set) ашиглах гэтэл техник хангамж reference бит өгдөггүй. Reference битийг хэрхэн симуляц хийх, эсвэл яагаад боломжгүйг тайлбарла. Боломжтой бол зардлыг тооцоол.}

\textbf{Боломжтой} — valid–invalid (эсвэл хамгаалалтын) битийг ашиглана. Үйлдлийн систем хуудас бүрт программын “reference бит” хадгална. Үе үе (жнь таймерын тасалдал бүрт) бүх хуудасны программын битийг 0 болгож, хуудасны хүснэгтийн бичлэгийг \emph{invalid} (эсвэл хандах эрхгүй) гэж тэмдэглэнэ — хуудас санах ойд хэвээр байна. Тэр хуудсанд дараагийн удаа хандахад trap гарч, OS хуудас санах ойд байгааг мэдээд программын reference битийг 1 болгож, бичлэгийг \emph{valid} болгоод зааврыг дахин эхлүүлнэ.

\textbf{Зардал:} үе бүрт ашиглагдсан хуудас бүрт нэг нэмэлт trap (хэдэн мкс, дискний I/O-гүй) ба бүх бичлэгийг цэвэрлэх ажил; мөн TLB-г цэвэрлэх шаардлагатай. Page fault-ийн дискний үйлдэлтэй (миллисекунд) харьцуулахад бага боловч үеийг хэт богино болговол trap-ийн ачаалал мэдэгдэхүйц болно.

%---------------------------------------------------------------
\question{10.11}{Та оновчтой байж магадгүй гэж бодож буй шинэ хуудас солих алгоритм зохиожээ. Зарим тусгай тохиолдолд Belady-ийн аномали гарч байна. Энэ алгоритм оновчтой юу? Тайлбарла.}

\textbf{Үгүй.} Оновчтой алгоритм (OPT) Belady-ийн аномалид өртдөггүй. Шинэ алгоритмыг $N$, $f(k)$-г $k$ фрэймтэй үеийн fault-ийн тоо гэе. Аномали гэдэг нь ямар нэг мөрөнд $f_N(k+1) > f_N(k)$ гэсэн үг. Харин OPT-ийн хувьд $f_{OPT}(k+1) \le f_{OPT}(k) \le f_N(k)$ тул $f_N(k+1) > f_{OPT}(k+1)$ — тэр мөр ба $k+1$ фрэйм дээр $N$ нь OPT-оос олон fault гаргаж байна. Тэгэхээр $N$ оновчтой биш.

%---------------------------------------------------------------
\question{10.12}{Сегментчлэл нь хуудаслалттай төстэй боловч хувьсах хэмжээтэй “хуудас” хэрэглэдэг. FIFO ба LRU хуудас солилтод суурилсан хоёр сегмент солих алгоритм тодорхойл. Солих гэж сонгосон сегмент шаардлагатай сегментэд хангалттай үргэлжилсэн зай чөлөөлөхгүй байж болохыг санаарай. Сегментийг шилжүүлж (relocate) болохгүй ба болох системийн стратегийг авч үз.}

\textbf{FIFO-д суурилсан:} сегментүүдийг санах ойд орсон дарааллаар нь дараалалд хадгална. Шинэ $s$ хэмжээтэй сегмент ирэхэд хангалттай нүх байхгүй бол дарааллын толгойн (хамгийн эрт орсон) сегментийг гаргана; үүсэх нүх $\ge s$ болтол давтана.

\textbf{LRU-д суурилсан:} сегмент бүрийн сүүлд хандсан хугацааг хадгалж, хамгийн удаан хэрэглэгдээгүй сегментийг эхэлж гаргана; хангалттай нүх гартал давтана.

\begin{itemize}[leftmargin=1.5em]
  \item \textbf{Шилжүүлэх боломжгүй систем:} нийт чөлөөт зай хангалттай байлаа ч \emph{үргэлжилсэн} байх ёстой. Тиймээс дарааллаар нь сохроор гаргахын оронд: гаргахад $s$-ээс том үргэлжилсэн нүх үүсгэх (хөрш нүхнүүдтэйгээ нийлэх) сегментүүдийн дунд FIFO/LRU-гаар хамгийн “хуучин”-ыг сонгоно; эсвэл санах ойн бүс бүрд түүнийг чөлөөлөхөд гаргах сегментүүдийн хамгийн хуучин/LRU утгыг тооцож, хамгийн бага хохиролтой үргэлжилсэн бүсийг сонгоно. Хэрэггүй сегментийг гаргахаас зайлсхийнэ.
  \item \textbf{Шилжүүлэх боломжтой систем:} FIFO/LRU дарааллаар сегмент гаргасаар нийт чөлөөт зай $\ge s$ болмогц санах ойг шахаж (compaction) бүх нүхийг нэгтгэнэ — шаардлагагүй сегментийг гаргахгүй ч шахалтын хуулбарлах зардалтай.
\end{itemize}

%---------------------------------------------------------------
\question{10.13}{Шаардлагаар хуудаслах системийн мультипрограммчлалын түвшин 4-т тогтоогдсон. CPU ба хуудаслах дискний ашиглалтыг хэмжсэн гурван үр дүн: (a) CPU 13\%, диск 97\%; (b) CPU 87\%, диск 3\%; (c) CPU 13\%, диск 3\%. Тохиолдол бүрт юу болж байна вэ? Мультипрограммчлалын түвшинг өсгөж CPU ашиглалтыг нэмэгдүүлж болох уу? Хуудаслалт тус болж байна уу?}

\begin{enumerate}[label=(\alph*), leftmargin=1.5em]
  \item \textbf{Thrashing.} Процессууд ихэнх хугацаагаа page fault-д зарцуулж, хуудаслах диск дүүрэн ачаалалтай, CPU хүлээж байна. Түвшинг өсгөвөл улам дордоно — \emph{бууруулах} (процесс түр гаргах) эсвэл санах ой нэмэх хэрэгтэй. Хуудаслалт тус болохгүй, харин хор болж байна.
  \item CPU өндөр ашиглалттай, page fault бага — систем хэвийн, үр дүнтэй ажиллаж байна. Түвшинг өсгөх шаардлага бараг байхгүй (CPU аль хэдийн 87\%), өсгөвөл thrashing руу ойртох эрсдэлтэй. Хуудаслалт тус болж байна (санах ойг хэмнэж байна).
  \item CPU ч, диск ч бага ашиглагдаж байна — ажиллах процесс хангалтгүй (эсвэл бусад I/O хүлээж байна). Мультипрограммчлалын түвшинг \textbf{өсгөж} CPU ашиглалтыг нэмэгдүүлж болно; хуудаслалт асуудал биш.
\end{enumerate}

%---------------------------------------------------------------
\question{10.14}{Base ба limit регистр ашигладаг машиныг хуудасны хүснэгттэй болгож өөрчилжээ. Хуудасны хүснэгтээр base ба limit регистрийг симуляц хийж болох уу? Хэрхэн, эсвэл яагаад болохгүй вэ?}

\textbf{Болно}, гэхдээ хуудасны хэмжээний нарийвчлалтай. Процессын санах ойг base хаягаас эхлэх \emph{дараалсан} фрэймүүдэд байрлуулж, логик хуудас $i$-г фрэйм $\text{base}/\text{хуудас} + i$ руу заана; $\lceil \text{limit}/\text{хуудас} \rceil$-ээс хойших бүх бичлэгийг \emph{invalid} болгоно. Ингэвэл хаяг хөрвүүлэлт нь “base нэмэх”-тэй, invalid бичлэг нь “limit-ээс хэтэрсэн бол trap”-тай ижил ажиллана.

Хязгаарлалт: base нь хуудасны хил дээр байх ёстой, limit нь хуудасны хэмжээ рүү дээш бөөрөнхийлөгдөнө — сүүлийн хуудасны үлдэгдэлд хандахад trap гарахгүй (дотоод хоосон зай). Нарийн нэг байтын нарийвчлалтай limit хэрэгтэй бол хуудасны хүснэгт дангаараа хангалтгүй.

%---------------------------------------------------------------
\section{Programming Problem 10.44 (бичлэгт)}

\texttt{lab9/page\_replace.c}: 0–9 хуудасны дугаартай санамсаргүй лавлагааны мөр үүсгэж, FIFO, LRU, OPT алгоритмыг ижил мөр дээр ажиллуулаад page fault-ийн тоог хэвлэнэ. Фрэймийн тоо командын мөрөөс; урт ба seed-ийг сонгож болно; \texttt{-r} сонголтоор тодорхой мөр өгч номын жишээгээр шалгана.
\begin{itemize}[leftmargin=1.5em]
  \item \textbf{FIFO} — тойрог дарааллын толгойн фрэймийг солино.
  \item \textbf{LRU} — фрэйм бүрт сүүлд хандсан индексийг хадгалж, хамгийн бага индекстэйг солино.
  \item \textbf{OPT} — фрэйм бүрийн хуудас ирээдүйд хэзээ хэрэглэгдэхийг хайж, хамгийн хол (эсвэл хэзээ ч үгүй)-г солино.
\end{itemize}
\begin{lstlisting}[language={}]
$ ./page_replace 3 -r "7 0 1 2 0 3 0 4 2 3 0 3 2 1 2 0 1 7 0 1"
Page frames: 3
Reference string (20): 7 0 1 2 0 3 0 4 2 3 0 3 2 1 2 0 1 7 0 1
FIFO: 15 page faults
LRU:  12 page faults
OPT:  9 page faults
$ ./page_replace 4 20 7          # санамсаргүй мөр, 4 фрэйм, seed 7
\end{lstlisting}
Номын §10.4.2–10.4.4-ийн жишээний хариу (FIFO 15, OPT 9, LRU 12)-тай таарч байна. 10.8, 10.9-ийн хүснэгтийг мөн энэ программаар давхар шалгасан.

\end{document}
